\documentstyle[%


righttag%


,11pt%


,amssymb%


,russian%


,izvru%


]{amsart}





\newcommand{\tl}[3]{\medskip\noindent{\bf#1}\ #2 \dotfill #3 \par} 


\pagestyle{empty}


\begin{document}


%\tableofcontents


\par


\vskip 2cm


\par


\bigskip





Номер посвящен исследованиям по теории дифференциальных и


функционально-диф\-фе\-рен\-ци\-аль\-ных уравнений. Изучаются краевые задачи


и задачи управления. Предыдущие номера по данной тематике:


\No\,5, 1993; \No\,6, 1994; \No\,10, 1995; \No\,11, 1996;


\No\,6, 1997; \No\,2, 1999; \No\,6, 2001;\linebreak \No\,6, 2002; \No\,4, 2003.





\bigskip


\bigskip


\par


\centerline{\Large СОДЕРЖАНИЕ}


\bigskip


\bigskip


%%%%%%%%%% Use \newline %%%%%%%%%%%











\tl{Абдурагимов Г.Э.} 


{О существовании положительного решения краевой задачи для одного\newline  


нелинейного функционально-дифференциального уравнения второго  


порядка}{3}


\tl{Бояринцев Ю.Е., Орлова И.В.} 


{Блочные алгебро-дифференциальные системы и их\newline индексы}{6}


\tl{Бравый Е.И.}


{О разрешимости двухточечной краевой задачи для линейного


сингуляр-\newline ного функционально-дифференциального уравнения}{14}


\tl{Власов В.В., Медведев Д.А.}


{Об оценках решений дифференциальных уравнений с


за-\newline паздывающим аргументом}{21}


\tl{Култышев С.Ю., Култышева Л.М.}


{Об идентификации некоторых классов


оператор-\newline ных моделей эволюционного типа}{30}


\tl{Мокейчев В.С.}


{Об отсутствии  собственных значений у


задачи, обобщающей задачу о\newline волноводах}{41}


\tl{Сапронов И.В.} 


{Об одном классе решений уравнения Вольтерра II рода с регулярной\newline 


особенностью в банаховом пространстве}{48}


\tl{Черепенников В.Б.}


{Обратная начальная задача для некоторых линейных систем


диффе-\newline ренциально-разностных уравнений}{59}


\tl{Чудинов К.М.} 


{Об устойчивости по части переменных линейных  


автономных систем с\newline последействием}{72}


\tl{Шрагин И.В.} 


{Классификация типа Бэра измеримых и стандартных функций}{81}

















\vfill


\noindent\raisebox{2pt}{\copyright} Казанский госудаpственный унивеpситет


\end{document}


